% =============================================================
% Annexe E, section 4 (ex-annexe F) — Jeu d'évaluation de l'enrichissement sémantique
% Alimentée depuis Ch.6 (§6.3.4).
% Condensation vague 3 : spécification du jeu et méthodes comparées ramenées
% en prose (deux tableaux supprimés) ; empreinte du catalogue ATT&CK renvoyée
% au ch.6 plutôt que redite ; porte de qualité int8 ramenée à un paragraphe.
% =============================================================
\section{Jeu d'évaluation de l'enrichissement sémantique}
\label{ann:evaluation-semantique}

Cette section donne le détail de l'évaluation condensée au \cref{ch:validation} :
constitution du jeu de test, méthodes comparées, grandeurs relevées et condition de publication de
chacune.

\medskip
\noindent\textbf{Jeu d'évaluation et méthodes comparées.} Le jeu est prélevé sur les alertes produites par le système durant une période d'observation
définie : \textbf{100 alertes visées, 50 au minimum}, tirées aléatoirement parmi les alertes
distinctes afin d'éviter la sur-représentation des règles bruyantes. Chaque alerte est rattachée
à la main à la technique attendue, ou à l'étiquette \og aucune technique applicable \fg{}, dans le
catalogue ATT\&CK Entreprise chargé dans l'entrepôt --- instantané non épinglé, identifié par
l'empreinte publiée en \cref{sec:eval-detection}. Une personne réalise l'annotation initiale et une
seconde réannote le sous-ensemble de contrôle : c'est cette seconde annotation qui conditionne la
publication.

Trois méthodes de rattachement sont mises en regard sur ce même jeu. \textbf{M0}, correspondance
par mots-clés adossée à une table écrite à la main entre termes fréquents et techniques, sert de
référence basse. \textbf{M1}, représentation statistique du texte et similarité cosinus contre les
descriptions de techniques, sert de référence intermédiaire, sans apprentissage de domaine.
\textbf{M2}, modèle de représentation spécialisé retenu par le projet
\cite{securebert2022,smet2023}, est la méthode évaluée.

Une annotation réalisée par une seule personne, également auteure du système évalué, introduit un
biais que le protocole traite plutôt que de l'ignorer (\cref{sec:limites}) : ce sont les deux
atténuations exigées --- annoter \emph{avant} de consulter les sorties du modèle, et faire réannoter
un sous-ensemble par une seconde personne --- et non le temps disponible, qui commandent la
publication d'un chiffre comparatif. Une exactitude au premier rang produite par un annotateur
unique confondu avec l'auteur du système ne mesurerait pas la qualité du rattachement, mais un
accord de l'auteur avec lui-même.

\medskip
\noindent\textbf{Grandeurs relevées et conditions de publication.} Quatre grandeurs de qualité sont définies avant tout relevé. L'\textbf{exactitude au premier rang}
est la part des alertes dont la technique attendue est proposée en première position ;
l'\textbf{exactitude aux trois premiers rangs} étend le compte aux trois propositions ; le
\textbf{taux d'abstention} est la part des alertes laissées sans rattachement, sous le seuil de
similarité ; le \textbf{taux d'erreur silencieuse} est la part des alertes rattachées à une
technique \emph{incorrecte} avec un score supérieur au seuil --- la métrique la plus importante et
la moins souvent rapportée. Une abstention est visible et sans danger, l'analyste sachant que le
système ne sait pas ; un rattachement erroné présenté avec un score élevé oriente vers la mauvaise
procédure. \textbf{Un système qui s'abstient souvent mais se trompe rarement est préférable à
l'inverse}, et cette préférence doit rester visible dans la mesure.

Le protocole ne publie aucune valeur avant que ses conditions de validité ne soient réunies, et ces
conditions sont écrites à l'avance plutôt que constatées après coup
(\cref{tab:resultats-methodes}). Les six grandeurs sont relevées pour les trois méthodes sur le
\emph{même} échantillon, en une seule exécution, faute de quoi la comparaison porterait sur des
jeux différents et ne mesurerait rien.

\begin{table}[htbp]
\caption{Grandeurs relevées et condition de publication de chacune}
\label{tab:resultats-methodes}
\centering
\scriptsize
\begin{tabular}{|P{3.7cm}|P{4.7cm}|P{6.9cm}|}
\hline
\textbf{Grandeur} & \textbf{Relevée pour} & \textbf{Condition de publication} \\
\hline
Exactitude au premier rang & M0, M1, M2 & Jeu de 50 à 100 alertes annotées, échantillonné
aléatoirement parmi les alertes distinctes \\ \hline
Exactitude aux trois premiers rangs & M1, M2 (sans objet pour M0, qui ne classe pas) & Idem \\ \hline
Taux d'abstention & M0, M1, M2 & Seuil de similarité déclaré et identique pour les trois
méthodes \\ \hline
Taux d'erreur silencieuse & M0, M1, M2 & Annotation réalisée \emph{avant} consultation des
sorties du modèle \\ \hline
Latence par alerte & M1, M2 & Relevée sur la même machine et le même lot \\ \hline
Coût pour mille alertes & M2 & Rapporté au dimensionnement déclaré du composant d'encodage \\ \hline
\textbf{Toutes} & --- & \textbf{Taux d'accord inter-annotateurs calculé sur un sous-ensemble
réannoté par une seconde personne} \\
\hline
\end{tabular}
\end{table}

La dernière ligne est la condition dominante : c'est la seule qui ne dépende pas de la machine mais
d'un tiers, d'où sa formulation comme exigence de méthode et non comme étape de calendrier. Trois
questions guident ensuite l'interprétation, quelle que soit la direction du résultat : le gain de
M2 sur M1 existe-t-il --- un résultat négatif est admis par avance au titre de l'objectif O5 et
serait rapporté tel quel ; ce gain justifie-t-il le coût d'un composant supplémentaire, que
mesurent la latence et le coût du \cref{tab:resultats-methodes} ; et sur quel type d'alerte se
concentre-t-il, une analyse par type étant plus informative qu'un chiffre global.

Une porte de qualité distincte a déjà été exécutée sur M2 lui-même, avant son emploi en recette :
la bascule vers un export quantifié (int8), comparée aux rapprochements du modèle de référence
(fp32) sur cinq requêtes de type SOC le 1\textsuperscript{er} août 2026, a rendu 0~correspondance
de premier rang sur 5, ce qui a fait rejeter cet export et conserver le format non quantifié au
prix d'une empreinte mémoire plus élevée (\cref{sec:eval-enrichissement}). Ce test ne se substitue
pas à la comparaison M0/M1/M2 : il montre que cet export int8 précis n'est pas acceptable sur le
petit échantillon retenu, il ne valide pas la qualité générale de fp32.

% --- Rétablissement des blancs par défaut autour des flottants : le formulaire de dépôt
%     qui suit n'appartient pas au périmètre de cette équipe. ---
\setlength{\textfloatsep}{20pt plus 2pt minus 4pt}
\setlength{\intextsep}{12pt plus 2pt minus 2pt}
\setlength{\floatsep}{12pt plus 2pt minus 2pt}
